
\section{Lower Bound: Achieving high accuracy requires abstaining}
%
Choosing thresholds for online anomaly detection is non-trivial.
In the absence of abstaining, the expected sum of FPs and FNs is at-least order $\sqrt{T}$ on an i.i.d. stream. We provide a sketch and defer details to the Appendix. Consider the standard reduction of estimation to hypothesis testing \cite{wainwright2019high} where the stream is either coming from the standard gaussian or a unit variance gaussian with mean located at $\frac{1}{\sqrt{T}}$. Denote by $Q^{(0)}(p) < Q^{(1)}(p)$ to be the $p$th quantile of the two distributions respectively. Lower bounds from hypothesis testing gives that at any test that looks at all $T$ samples and identifies which of the two distributions the stream came from makes a mistake with probability at-least $1/8$. Furthermore, Chernoff bound applied to the binary random variables $\mathbf{1}_{Q^{(0)}(p) < S_t < Q^{(1)}(p)}$ gives that with probability at-least $1-e^{\sqrt{T}/8}$, at-least ${C_p\sqrt{T}}$ samples among $S_1, \cdots, S_T$ lies in the range $[Q^{(0)}(p), Q^{(1)}(p)]$, where $C_p$ is a constant depending on $p$. Thus, an union bound gives that in the absence of abstains, with probability at-least $1/8-e^{\sqrt{T}/8}$, the sum of FP and FN is at-least ${C_p\sqrt{T}}$. A formal statement and proof in Section \ref{sec:lower_bound_proof} in the Appendix.
